% chapters/06/section-04.tex
\section{Pantulan dengan Koefisien Restitusi}
\label{sec:pantulan}

Benda yang menyentuh lantai tidak selalu harus berhenti. Bagian ini membahas
pantulan, yaitu pembalikan arah gerak yang disertai pengurangan kecepatan.

\subsection{Pembalikan Arah dan Pengurangan Kecepatan}

Pantulan terjadi dalam dua langkah: arah gerak dibalik, lalu besarnya
dikurangi. Keduanya dilakukan sekaligus dengan satu perhitungan.

\begin{kodekeluaran}
velocity.y = -velocity.y * restitution
\end{kodekeluaran}

Perhatikan tanda negatif dan pengalinya. Tanda negatif membalik arah, sedangkan
pengali \code{restitution} mengurangi besarnya kecepatan. Nilai pengali itu
terletak antara nol dan satu, dengan makna sebagai berikut.

\begin{itemize}
  \item \textbf{Nilai 0} berarti tidak ada pantulan sama sekali: benda
        berhenti seketika. Nilai ini dipakai untuk pemain, sebagaimana
        dijelaskan pada Bagian~\ref{sec:gravitasi-lantai}.
  \item \textbf{Nilai 1} berarti pantulan sempurna: kecepatan tidak
        berkurang, dan benda memantul selamanya. Nilai ini jarang dipakai
        karena membuat permainan tidak pernah tenang.
  \item \textbf{Nilai antara 0 dan 1}, misalnya 0,8, berarti benda kehilangan
        sebagian kecepatannya setiap pantulan sehingga akhirnya berhenti.
\end{itemize}

\subsection{Tinggi Pantulan Menyusut Sebesar Kuadrat Koefisien}

Kecepatan setelah pantulan adalah $e$ kali kecepatan sebelumnya. Tinggi yang
dapat dicapai sebanding dengan \textbf{kuadrat} kecepatan, sehingga tinggi
setelah pantulan tinggal $e^2$ kali tinggi sebelumnya. Untuk $e = 0{,}8$,
faktornya adalah $0{,}64$: setiap pantulan menyisakan 64 persen tinggi
sebelumnya. Susutan itu diperlihatkan pada Gambar~\ref{fig:pantulan}.

\begin{figure}[htbp]
\centering
% chapters/06/figures/pantulan.tikz
% Pantulan dengan koefisien restitusi 0,8: setiap pantulan menyisakan 64
% persen tinggi sebelumnya, sebab tinggi sebanding dengan kuadrat kecepatan.
\begin{tikzpicture}[
  bola/.style={circle, draw=ubAccent, thick, fill=ubAlt, inner sep=2pt},
  lantai/.style={thick, draw=ubLine},
  label/.style={font=\scriptsize, align=left},
  lintasan/.style={very thick, draw=ubPrimary},
  ukur/.style={draw=ubGreen, dashed, thin}
]

\draw[lantai] (0,0) -- (8.2,0);

% Jatuh dari ketinggian h.
\draw[lintasan] (0,3.2) .. controls (0.75,1.5) and (1.0,0.5) .. (1.2,0);

% Tiga pantulan berikutnya, tinggi menyusut menjadi 0,64 kali setiap kali.
\draw[lintasan] (1.2,0) .. controls (1.75,1.7) and (2.25,2.05) .. (2.6,2.05);
\draw[lintasan] (2.6,2.05) .. controls (2.95,2.05) and (3.5,1.6) .. (4.0,0);

\draw[lintasan] (4.0,0) .. controls (4.35,1.1) and (4.9,1.31) .. (5.1,1.31);
\draw[lintasan] (5.1,1.31) .. controls (5.35,1.31) and (5.75,1.0) .. (6.2,0);

\draw[lintasan] (6.2,0) .. controls (6.4,0.7) and (6.8,0.84) .. (6.9,0.84);
\draw[lintasan] (6.9,0.84) .. controls (7.05,0.84) and (7.3,0.6) .. (7.7,0);

% Bola pada puncak setiap pantulan.
\node[bola] at (0,3.2) {};
\node[bola] at (2.6,2.05) {};
\node[bola] at (5.1,1.31) {};
\node[bola] at (6.9,0.84) {};

% Garis ukur tinggi setiap pantulan.
\draw[ukur] (0,3.2) -- (8.6,3.2);
\draw[ukur] (0,2.05) -- (8.6,2.05);
\draw[ukur] (0,1.31) -- (8.6,1.31);
\draw[ukur] (0,0.84) -- (8.6,0.84);

\node[label, ubGreen, anchor=west] at (8.7,3.2) {$h$};
\node[label, ubGreen, anchor=west] at (8.7,2.05) {$0{,}64\,h$};
\node[label, ubGreen, anchor=west] at (8.7,1.31) {$0{,}41\,h$};
\node[label, ubGreen, anchor=west] at (8.7,0.84) {$0{,}26\,h$};

\node[label, anchor=north west, text width=150mm] at (0,-0.35)
  {Dengan koefisien restitusi $e = 0{,}8$, kecepatan setelah setiap pantulan
   tinggal $0{,}8$ kali, sedangkan tinggi yang dicapai tinggal
   $e^2 = 0{,}64$ kali, sebab tinggi sebanding dengan kuadrat kecepatan.};
\end{tikzpicture}

\caption{Tinggi pantulan yang menyusut dengan koefisien restitusi 0,8}
\label{fig:pantulan}
\end{figure}

Kuadrat itulah sebabnya pantulan terasa cepat mereda walaupun kecepatannya
hanya berkurang dua puluh persen setiap kali. Pada pantulan ketiga, tinggi yang
dicapai hanya tinggal seperempat tinggi semula.

\subsection{Urutan Penanganan yang Benar}

Satu hal yang mudah terlewat adalah urutan pekerjaan ketika benda menyentuh
lantai. Posisi harus dikembalikan ke permukaan lantai \textbf{lebih dahulu},
baru kecepatan diubah.

\begin{kodekeluaran}
if position.y + half >= floor_y:
    position.y = floor_y - half                   # 1. perbaiki posisi
    if velocity.y > 0:
        velocity.y = -velocity.y * restitution    # 2. balik kecepatan
\end{kodekeluaran}

Apabila urutannya dibalik, benda membalik arahnya sementara posisinya masih
terbenam di dalam lantai. Pada bingkai berikutnya, benda bergerak ke atas dari
posisi yang salah itu, dan sebagian kecepatannya terpakai untuk keluar dari
lantai. Gejalanya halus tetapi khas: tinggi pantulan menyusut lebih cepat
daripada yang seharusnya, dan pada benda yang cepat, benda dapat tersangkut di
dalam lantai.

Perhatikan juga syarat \code{if velocity.y > 0}. Syarat itu memastikan
pantulan hanya dilakukan oleh benda yang sedang bergerak ke bawah. Tanpa syarat
itu, benda yang baru saja memantul dan sedang bergerak ke atas akan dibalik
arahnya lagi pada bingkai berikutnya, sehingga benda terperangkap di lantai.

\subsection{Pantulan pada Dinding}

Pantulan pada dinding bekerja dengan cara yang sama, hanya sumbunya yang
berbeda: yang dibalik adalah komponen mendatar.

\begin{kodekeluaran}
if position.x - half < 0:
    position.x = half
    velocity.x = abs(velocity.x) * restitution
elif position.x + half > width:
    position.x = width - half
    velocity.x = -abs(velocity.x) * restitution
\end{kodekeluaran}

Perhatikan pemakaian \code{abs()} untuk menentukan arah sesudah pantulan.
Dinding kiri selalu mengarahkan benda ke kanan, dan dinding kanan selalu
mengarahkannya ke kiri, terlepas dari arah datangnya. Memakai tanda negatif
sederhana tidak cukup, karena benda yang sudah bergerak menjauhi dinding akan
dibalik kembali ke arah dinding dan terperangkap di sana.

\subsection{Menghentikan Pantulan yang Terlalu Kecil}

Benda yang memantul dengan koefisien restitusi akan terus memantul, walaupun
tingginya makin lama makin kecil. Pada akhirnya pantulan itu terlalu pendek
untuk terlihat, tetapi perhitungannya tetap berjalan dan membuat benda
bergetar di permukaan lantai, suatu gejala yang disebut \istilah{jitter}.

Penanganannya sederhana: hentikan pantulan apabila kecepatannya sudah di bawah
nilai tertentu.

\begin{kodekeluaran}
velocity.y = -velocity.y * restitution
if abs(velocity.y) < MIN_BOUNCE_SPEED:
    velocity.y = 0.0
\end{kodekeluaran}

Nilai batas pada program bab ini adalah 60 piksel per detik. Nilai itu dipilih
karena pantulan dengan kecepatan lebih kecil daripada itu hanya mengangkat
benda setinggi satu piksel, sehingga tidak terlihat tetapi tetap memakan
perhitungan.

\begin{tipsbox}
Cara yang sama dipakai untuk pemain, yang koefisien
restitusinya nol. Pada pemain, kecepatan tegaknya langsung menjadi nol tanpa
memerlukan pemeriksaan batas, sebab hasil perkalian dengan nol memang bernilai
nol. Satu baris kode itu melayani dua keperluan sekaligus, dan itulah manfaat
menuliskan hukum fisika secara umum.
\end{tipsbox}
